% Problem record op_36ac6718d2c37628. This TeX file is derived from
% database/problems_json/op_36ac6718d2c37628.json by scripts/sync-tex.mjs; edit the JSON record
% and rerun the script rather than editing this file.
\section{Polynomial-time quantum algorithm for approximate Shortest Vector Problem}
\label{sec:op_36ac6718d2c37628}

\paragraph{Problem.}
Does the polynomial-factor approximate Shortest Vector Problem admit a
polynomial-time quantum algorithm?

Let $B\in\mathbb Q^{n\times n}$ be a nonsingular lattice basis whose entries
have polynomially bounded bit length, and let
\begin{equation}
  \mathcal L(B)=\{Bz:z\in\mathbb Z^n\}\subset\mathbb R^n .
  \label{eq:36ac-appsvp-lattice}
\end{equation}
Define the length of a shortest nonzero lattice vector by
\begin{equation}
  \lambda_1(\mathcal L)=\min_{v\in\mathcal L\setminus\{0\}}|v|_2 .
  \label{eq:36ac-appsvp-minimum}
\end{equation}
For an approximation factor $\gamma=\gamma(n)\geq 1$, the $\gamma$-approximate
Shortest Vector Problem asks for a nonzero vector $v\in\mathcal L(B)$
satisfying
\begin{equation}
  |v|_2\leq\gamma(n)\lambda_1(\mathcal L(B)).
  \label{eq:36ac-appsvp-goal}
\end{equation}
The open question is whether, for polynomial approximation factors such as
$\gamma(n)=n^c$ for a fixed constant $c>0$, there exists a bounded-error
quantum algorithm that outputs a vector satisfying
\eqref{eq:36ac-appsvp-goal} in time polynomial in the bit length of $B$.

\subsection*{Status}

Unsolved

\subsection*{Source}

Contributor: unknown.  The opening question is a contributor-formulated
synthesis of a standard open algorithmic problem in the lattice and
quantum-algorithms literature.  Biasse, Bonnetain, Kirshanova,
Schrottenloher, and Song explicitly identify polynomial-factor approximate
SVP as a principal regime of interest and survey the known quantum
algorithms
\sourcecite{ref:36ac-biasse}{Bia23}, but do not state the opening sentence
verbatim as a named open question.  It should not be attributed verbatim to
any of the cited authors.

\subsection*{Progress}

\begin{itemize}
  \item Classical polynomial-time lattice reduction does not reach the
polynomial-factor regime in general: the LLL algorithm gives only an
approximation factor exponential in $n$ for arbitrary lattices, so a
polynomial approximation factor in polynomial time would be a qualitatively
stronger result rather than an improvement of the approximation constant in
LLL
\sourcecite{ref:36ac-lll}{LLL82}.

  \item Known algorithms obtain substantially better approximation factors by
spending exponential time: Aggarwal, Li, and Stephens-Davidowitz gave a
$2^{n/2+o(n)}$-time algorithm achieving an approximation factor
$\widetilde O(\sqrt{n})$ for both SVP and Hermite SVP (the tilde
suppresses polylogarithmic factors), together with broader
time--approximation tradeoffs.  Polynomial approximation factors are
therefore accessible significantly faster than exact SVP in some regimes,
but not in polynomial time
\sourcecite{ref:36ac-als}{ALS21}.

  \item Quantum computation provides speedups for exponential-time SVP
approaches but no known technique reaches polynomial time: quantum
lattice-sieving algorithms use amplitude amplification and quantum walks to
reduce the exponential constant, and Bonnetain, Chailloux, Schrottenloher,
and Shen obtained heuristic quantum sieving complexity
$2^{0.2563n+o(n)}$ for SVP
\sourcecite{ref:36ac-bcss}{BCSS23}.

  \item There are also provable quantum algorithms outside the heuristic-sieving
framework: Aggarwal, Chen, Kumar, and Shen gave a provable quantum algorithm
for exact SVP running in $2^{0.950n+o(n)}$ time with exponential classical
memory and polynomially many qubits, and in the QRAM model in
$2^{0.835n+o(n)}$ time; since an exact-SVP solver also solves the
approximate problem, these are quantum upper bounds for approximate SVP but
remain exponential
\sourcecite{ref:36ac-acks}{ACKS25}.

  \item The problem is distinct from the unique Shortest Vector Problem: Regev
showed that an efficient solution to the relevant dihedral coset problem
--- in particular an appropriate coset-sampling solution of the Dihedral
Hidden Subgroup Problem --- would yield a polynomial-time quantum algorithm
for polynomial-factor unique-SVP.  The uniqueness promise is absent in the
arbitrary-lattice problem here, so the DHSP--lattice reduction does not give
a polynomial-time algorithm for general approximate SVP
\sourcecite{ref:36ac-regev04}{Reg04}.

  \item Approximate SVP is likewise distinct from GapSVP and SIVP: Regev's
worst-case-to-average-case reduction for Learning With Errors is a quantum
hardness reduction from worst-case approximate GapSVP and SIVP to LWE, not
an algorithm for the search problem here.  The LWE, DHSP, and
approximate-SVP questions are closely related but remain logically distinct
open problems
\sourcecite{ref:36ac-regev09}{Reg09}.

  \item A survey of quantum cryptanalysis by Biasse, Bonnetain, Kirshanova,
Schrottenloher, and Song treats polynomial-factor approximate SVP as the
principal approximation regime relevant to lattice-based cryptographic
constructions and surveys the known quantum enumeration, sieving, and
BDD-based approaches; these methods provide at most exponential-time
speedups in the general setting
\sourcecite{ref:36ac-biasse}{Bia23}.
\end{itemize}

\subsection*{References}

\begin{enumerate}[label={},leftmargin=4.8em,labelsep=0.5em]
  \footnotesize
  \item[\textup{[LLL82]}]\label{ref:36ac-lll}
  Arjen K. Lenstra, Hendrik W. Lenstra, Jr., and L\'aszl\'o Lov\'asz,
``Factoring Polynomials with Rational Coefficients,'' \emph{Mathematische
Annalen} \textbf{261}, 515--534 (1982).
\href{https://doi.org/10.1007/BF01457454}{doi:10.1007/BF01457454}.

  \item[\textup{[Reg04]}]\label{ref:36ac-regev04}
  Oded Regev, ``Quantum Computation and Lattice Problems,'' \emph{SIAM
Journal on Computing} \textbf{33}, 738--760 (2004).
\href{https://doi.org/10.1137/S0097539703440678}{doi:10.1137/S0097539703440678};
\href{https://arxiv.org/abs/cs/0304005}{arXiv:cs/0304005}.

  \item[\textup{[Reg09]}]\label{ref:36ac-regev09}
  Oded Regev, ``On Lattices, Learning with Errors, Random Linear Codes,
and Cryptography,'' \emph{Journal of the ACM} \textbf{56}, Article 34
(2009).
\href{https://doi.org/10.1145/1568318.1568324}{doi:10.1145/1568318.1568324}.

  \item[\textup{[ALS21]}]\label{ref:36ac-als}
  Divesh Aggarwal, Zeyong Li, and Noah Stephens-Davidowitz, ``A
$2^{n/2}$-Time Algorithm for $\sqrt{n}$-SVP and $\sqrt{n}$-Hermite SVP, and
an Improved Time-Approximation Tradeoff for (H)SVP,'' in \emph{Advances in
Cryptology -- EUROCRYPT 2021}, Part I, 467--497 (2021).
\href{https://doi.org/10.1007/978-3-030-77870-5_17}{doi:10.1007/978-3-030-77870-5_17};
\href{https://arxiv.org/abs/2007.09556}{arXiv:2007.09556}.

  \item[\textup{[BCSS23]}]\label{ref:36ac-bcss}
  Xavier Bonnetain, Andr\'e Chailloux, Andr\'e Schrottenloher, and Yixin
Shen, ``Finding Many Collisions via Reusable Quantum Walks -- Application
to Lattice Sieving,'' in \emph{Advances in Cryptology -- EUROCRYPT 2023},
Part V, 221--251 (2023).
\href{https://doi.org/10.1007/978-3-031-30589-4_8}{doi:10.1007/978-3-031-30589-4_8};
\href{https://arxiv.org/abs/2205.14023}{arXiv:2205.14023}.

  \item[\textup{[ACKS25]}]\label{ref:36ac-acks}
  Divesh Aggarwal, Yanlin Chen, Rajendra Kumar, and Yixin Shen,
``Improved Classical and Quantum Algorithms for the Shortest Vector Problem
via Bounded Distance Decoding,'' \emph{SIAM Journal on Computing}
\textbf{54}, 233--278 (2025).
\href{https://doi.org/10.1137/22M1486959}{doi:10.1137/22M1486959};
\href{https://arxiv.org/abs/2002.07955}{arXiv:2002.07955}.

  \item[\textup{[Bia23]}]\label{ref:36ac-biasse}
  Jean-Fran\c{c}ois Biasse, Xavier Bonnetain, Elena Kirshanova, Andr\'e
Schrottenloher, and Fang Song, ``Quantum Algorithms for Attacking Hardness
Assumptions in Classical and Post-Quantum Cryptography,'' \emph{IET
Information Security} \textbf{17}, 171--209 (2023).
\href{https://doi.org/10.1049/ise2.12081}{doi:10.1049/ise2.12081}.
\end{enumerate}

\subsection*{Comment}

The open question is whether any quantum algorithm breaks the
polynomial-time barrier for polynomial approximation factors.  Known quantum
algorithms for SVP and its variants remain exponential in the lattice
dimension, and the LWE and DHSP reductions provide hardness reductions or
require structural promises (uniqueness of the shortest vector) absent here.
See the related problems on \href{https://qiqc-op.com/problem/op_3596f8c955c66d94/}{Learning With Errors} and \href{https://qiqc-op.com/problem/op_90a05e57e44c086e/}{the Dihedral Hidden Subgroup Problem}.

\subsection*{Field}

Quantum algorithm

\subsection*{Topic}

Computational complexity and computability

\subsection*{ID}

\texttt{op\_36ac6718d2c37628}
